\subsection{与扭曲微分形式相关的测度}

回顾一下，在本节中，我们假定 K = R，且所考虑的所有流形都是局部有限维的。

10.4.1. 设 X 是类 \(C^{r}\) 的流形。取切丛 T(X) 作为向量丛 M，便可应用 10.3 的定义和结果。于是有 \(\mathrm{Or}_{\mathbf{M}} = \tilde{\mathbf{X}}\)（10.2.4）；我们将其记作 \(\tilde{\mathbf{R}}_{\mathbf{X}}\)，并将扭曲标量丛简单地称为丛 \(\tilde{\mathbf{R}}_{\mathrm{T}(\mathbf{X})}\)（10.3.1）\footnote{注意，当把 X 看作自身上秩为 0 的向量丛结构时，这一记号与 \(\widetilde{\mathbf{R}}_{\mathbf{M}}\) 的记号冲突。}；取值于向量丛 F 的 T(X)-扭曲微分形式简单地称为取值于 F 的扭曲（或奇）微分形式。当 X 配备一个定向时，映射 \(\omega \mapsto \xi \otimes \omega\) 使得可以将通常微分形式（有时称为“偶”微分形式）与扭曲形式等同。

10.4.2. 设 \(f\colon\mathbf{X}\to\mathbf{Y}\) 是类 \(\mathbf{C}^{r}\) 流形的态射，且 \(\tilde{f}\colon\tilde{\mathbf{X}}\to\tilde{\mathbf{Y}}\) 是 \(f\) 的一个定向（10.2.5）。存在唯一一个同构 \(j\)，从 \(f^{*}(\widetilde{\mathbf{R}}_{\mathbf{Y}})\) 到 \(\widetilde{\mathbf{R}}_{\mathbf{X}}\)，使得 \(\tilde{x}=j(\pi(\tilde{x}),\tilde{f}(\tilde{x}))\) 对所有 \(\tilde{x}\in\tilde{\mathbf{X}}\) 成立。我们通过 \(j\) 将这两个丛等同。若 \(\omega\) 是 Y 上取值于向量丛 F 的扭曲形式，次数为 \(p\)，则逆像 \(f^{*}(\omega)\) 可与 X 上次数为 \(p\)、取值于 \(f^{*}(\mathrm{F})\) 的扭曲微分形式等同。若将 \(\tilde{\omega}\)（相应地 \(\widetilde{f^{*}(\omega)}\)）记作分别位于 \(\tilde{\mathbf{Y}}\)（相应地 \(\tilde{\mathbf{X}}\)）上、与 \(\omega\)（相应地 \(f^{*}(\omega)\)）对应的微分形式，如 10.3.3 中那样，则有 \(\tilde{f}^{*}(\tilde{\omega})=\widetilde{f^{*}(\omega)}\)。

当 F 是由 Banach 空间定义的平凡丛时，运算 \(f^*\) 与外微分交换（10.3.4）。

10.4.3. 设 X 是维数为 n 的纯流形，\(\omega\) 是 X 上取值于 Banach 空间 E 的 n 次扭曲微分形式。令 \(c = (\mathrm{U}, \varphi, \mathbf{R}^{n})\) 为 X 的一个图，令 \(\xi^{n}\) 为 \(R^{n}\) 的典范定向；以 \(u^{1}, \ldots, u^{n}\) 表示 \(R^{n}\) 上的坐标函数。存在唯一一个函数 \(f_{c}\)，定义在 \(\varphi(\mathrm{U})\) 上，取值于 E，使得

\[
\omega | \mathrm{U} = \varphi^ {*} (\xi^ {n} \otimes f _ {c}. d u ^ {1} \wedge \dots \wedge d u ^ {n}).
\]

称 \(\omega\) 是局部可积的，如果对于 X 的每个图 \(c = (\mathrm{U},\varphi ,\mathbb{R}^n)\)，相应函数 \(f_{c}\) 关于 \(\lambda_{\varphi (\mathrm{U})}^{\otimes n}\) 局部可积（其中 \(\lambda\) 表示 R 上的勒贝格测度）。假定这一点成立且 X 是分离的。在 X 上存在一个取值于 E 的向量测度 \(\alpha (\omega)\)，并且唯一地具有如下性质：对于 X 的每个图 \(c = (\mathrm{U},\varphi ,\mathbb{R}^n)\)，其在 \(\varphi\) 下的像，即测度 \(\alpha (\omega)\) 在 U 上的限制，是测度 \(f_{c}\lambda_{\varphi(\mathrm{U})}^{\otimes n}\)（INT, VI, § 2, no. 4）。称 \(\alpha (\omega)\) 为由 \(\omega\) 定义的测度，并且通常仍简记为 \(\omega\)。\(\alpha (\omega)\) 的支集包含于 \(\operatorname{Supp}\omega\)（\(\omega\) 的支集）中（把微分形式 \(\omega\) 的支集定义为满足 \(x\in X\) 且 \(\omega (x)\neq 0\) 的 x 的集合的闭包）。

若 \(f\) 是 X 上关于测度 \(\alpha(\omega)\) 局部可积的实函数，则微分形式 \(f\omega\) 局部可积，且有 \(\alpha(f\omega) = f.\alpha(\omega)\)。

若 \(g\) 是 X 上关于 \(\alpha(\omega)\) 本质可积的实函数，则其积分通常记作 \(\int_{\mathbf{X}} g\omega\)、\(\int g\omega\) 或 \(\int g(x)\omega(x)\)；人们避免将其记作 \(\int g d\omega\)，因为这可能与 \(d\omega\)（\(\omega\) 的外微分）混淆；当 \(\omega\) 属于类 \(C^{1}\) 时，该外微分有定义（并且还等于 0）。若 A 是 X 的一个子集，且其示性函数 \(\varphi_{\mathrm{A}}\) 关于 \(\alpha(\omega)\) 本质可积，则 \(\varphi_{\mathrm{A}}\) 的积分记作 \(\int_{A} \omega\)。若常数 1 关于 \(\alpha(\omega)\) 本质可积，则称 \(\omega\) 是可积的。

现在令 \(p\) 为整数，满足 \(0 \leqslant p \leqslant n\)，令 \(\omega\) 为 X 上取值于 Banach 空间 E、次数为 \(p\) 的扭曲微分形式。此外，令 Y 为 X 的一个维数为 \(p\) 的纯子流形；假设典范嵌入 \(i: Y \to X\) 配备了一个定向。于是，其逆像 \(i^{*}(\omega)\)（10.4.2）有时记作 \(\omega|Y\)。若 \(\omega|Y\) 可积，置

\[
\int_ {\mathbf {Y}} \omega = \int_ {\mathbf {Y}} \omega | \mathbf {Y}.
\]

10.4.4. 设 X 是一个配备定向 \(\xi\) 的 n 维纯流形。通常微分形式与扭曲微分形式之间的等同 \(\omega \mapsto \xi \otimes \omega\) 使得可以将前述结果应用于 X 上的 n 次形式；因此，每个局部可积的 n 次形式 \(\omega\) 都关联着 X 上的一个测度，仍记作 \(\omega\)。若 E = R，则这样的形式 \(\omega\) 局部具有可积模（10.1.5），与其对应的测度 \(\text{mod}(\omega)_{\mu}\)（10.1.6）恰好是 \(|\omega|\)（测度 \(\omega\) 的绝对值）（INT, III, § 1, no. 6）。

10.4.5. 例。--- 设 X 是一个定向的、分离的 1 维纯实流形，且 \(z: \mathrm{X} \to \mathbf{C}\) 是从 X 到 C 的可微映射。复微分形式 \(dz\) 在 X 上定义了一个也记作 \(dz\) 的复测度。这尤其适用于 X 是 C 的一个 1 维纯可微子流形、\(z\) 是 X 到 C 的嵌入的情形。

更具体地，设 X 是中心为 0、半径 \(\rho > 0\) 的圆；按照 10.2.8 b) 中的指示，并利用 C 与 \(\mathbf{R}^2\) 的通常等同，为 X 赋予定向。若 \(x \in \mathbf{X}\)，切空间 \(\mathrm{T}_x(\mathbf{X})\) 可与 \(\mathbf{R}ix\) 等同，后者是 \(\mathrm{T}_x(\mathbf{C}) = \mathbf{C}\) 中的直线，而所选定向是半直线 \(\mathbf{R}_+ix\)。若 \(f\) 是 X 上取值于复 Banach 空间的连续函数，则 \(f\) 关于测度 \(dz\) 的积分由公式给出

\[
\int_ {\mathbb {X}} f (z) d z = i \rho \int_ {0} ^ {2 \pi} f (\rho e ^ {i \alpha}) e ^ {i \alpha} d \alpha .
\]

例如，若 \(n\) 是一个整数，则有

\[
\int_ {\mathrm{X}} z ^ {n} d z = i \rho^ {n + 1} \int_ {0} ^ {2 \pi} e ^ {(n + 1) i \alpha} d \alpha = \left\{ \begin{array}{l l} 0 & \quad \text {if } n \neq - 1 \\ 2 i \pi & \quad \text {if } n = - 1. \end{array} \right.
\]

10.4.6（“复数情形”）。设 \(X^{c}\) 是一个分离的、纯维数为 \(m\) 的复解析流形，且 \(\omega\) 是次数为 \(m\) 的全纯微分形式，定义在 \(X^{c}\) 上。设 X 是 \(X^{c}\) 的底层实解析流形；形式 \(\omega\) 与 X 上的 \((m,0)\) 型形式等同（8.8.9）；令 \(\overline{\omega}\) 为其共轭（8.8.2）。形式 \(i^{m^{2}}\omega \wedge \overline{\omega}\) 是 X 上次数为 \(2m\) 的实微分形式；考虑到 X 的典范定向（10.2.7），此形式与 X 上的一个测度等同，而该测度正是 10.1.6 中定义的正测度 \(\operatorname{mod}(\omega)_{\mu}\)，其中 \(\mu\) 是由微分形式 \(i \, dz \wedge \overline{dz}\) 定义的测度，换言之，是测度 \(\lambda^{\otimes 2}\) 在 \(\mathbf{C}\) 上的两倍（按通常方式与 \(\mathbf{R}^{2}\) 等同）。

令 H 为由全纯形式 \(\omega\)（次数为 \(m\)，定义在 \(\mathbf{X}^c\) 上）组成的空间，使得测度 \(\omega \wedge \overline{\omega}\) 有界。对于 H 中的 \(\alpha, \beta\)，复测度 \(\alpha \wedge \overline{\beta}\) 有界；令

\[
(\alpha | \beta) = i ^ {m ^ {2}} \int_ {\mathbf {X}} \alpha \wedge \bar {\beta},
\]

从而得到 H 上的一个厄米形式，它使 H 成为 Hilbert 空间。

